\section{Introduction}

\subsection{Regulatory Context}

The ICH M7(R2) guideline recommends the use of two complementary in~silico approaches for assessing the mutagenic potential of pharmaceutical impurities: a rule-based expert system and a statistical QSAR model.\cite{ICH_M7_2017} This regulatory framework has driven considerable interest in developing predictive models for the Ames bacterial reverse mutation assay,\cite{Ames_1973,Mortelmans_2000} with extensions to other genotoxicity endpoints including in~vitro chromosome aberration (CA) and in~vivo micronucleus (MN) tests.\cite{OECD_473,OECD_474} Reliable QSAR models for these endpoints could reduce animal testing, accelerate safety assessments, and support the 3Rs principle.\cite{Russell_Burch_1959}

\subsection{The Overestimation Problem}

Published Ames QSAR models report a wide range of performance metrics, with accuracies typically spanning 0.80--0.90 and MCC values of 0.60--0.80.\cite{Hansen_2009,Xu_2012,Honma_2019,Li_2021} However, these numbers arise from different data splits, datasets, preprocessing pipelines, and evaluation protocols, making direct comparison difficult.

Several factors can inflate apparent model performance. First, random data splitting allows structurally similar compounds (sharing the same Murcko scaffold\cite{Bemis_Murcko_1996}) to appear in both training and test sets, creating an optimistic estimate of generalization to novel structures. Scaffold-based splitting has been recommended as a more realistic evaluation strategy,\cite{Sheridan_2013,Wu_MoleculeNet_2018} but the magnitude of performance difference between random and scaffold splits has not been systematically quantified across multiple genotoxicity endpoints.

Second, most large Ames datasets aggregate compounds from distinct sources---typically pharmaceutical compound libraries and industrial chemical registries.\cite{Hansen_2009,Honma_2019} These sources differ in composition and mutagenicity prevalence. When these source domains are mixed in training and test sets, a model may partly learn source-domain membership rather than toxicity mechanisms---a form of confounding that has been noted\cite{Tropsha_2010,Hanser_2016} but, to our knowledge, never systematically measured.

Third, preprocessing choices such as salt stripping, metal compound removal, and charge normalization can substantially alter molecular representations,\cite{Fourches_2010,Fourches_2016} yet these steps are often undocumented or applied uniformly without assessing endpoint-specific effects.

\subsection{Gaps in the Current Literature}

Most published studies evaluate a single endpoint using a single split method and focus their analysis on algorithm comparison.\cite{Yang_2019,Mayr_2016} Cross-domain generalization is rarely tested. Statistical significance of model differences is seldom reported.\cite{Demsar_2006} The relative contributions of split method, source-domain composition, and preprocessing choices to apparent model performance have not been disentangled within a single consistent experimental framework.

\subsection{Contributions}

This work makes the following contributions:

\begin{enumerate}
    \item We introduce a three-level evaluation framework that progressively increases evaluation rigor---in-domain scaffold split, random-versus-scaffold comparison, and cross-source-domain leave-domain-out---applied to a single locked model configuration for fair comparison.
    
    \item We demonstrate that, for the Ames mutagenicity benchmark studied here, source-domain confounding arising from the mixing of drug-like and industrial compounds with a 19.5-fold prevalence difference drove a scaffold-to-LDO MCC drop of 0.390, substantially exceeding the split-method effect ($\Delta = 0.046$).
    
    \item We show that while in-domain algorithm differences were small across all seven algorithms (MCC range 0.063 under a fixed feature/preprocessing configuration), cross-domain robustness varied meaningfully---Random Forest maintained MCC = 0.325 under domain shift where XGBoost collapsed to 0.122.
    
    \item We demonstrate that in~vivo micronucleus models showed some ranking signal (ROC-AUC = 0.860, PR-AUC = 0.237) after salt stripping but lacked reliable classification utility, with MCC confidence intervals crossing zero.
    
    \item We demonstrate endpoint-specific preprocessing effects, with salt stripping improving in~vivo best-model MCC by +0.22 while having mixed or negative effects on other endpoints.
    
    \item We provide an open-source pipeline with all experimental configurations, split seeds, and locked results tables for full reproducibility.
\end{enumerate}
